% !TEX program = xelatex
% ============================================================
%  خطاب تجاري — Business Letter Template
%  قالب عربي (RTL) لخطاب تجاري رسمي وفق بنية DIN 5008
%  مُكيّفة للاتجاه من اليمين إلى اليسار.
%  Compile with: xelatex main.tex (run twice)
% ============================================================
%  Features:
%    - Right-to-left (RTL) layout with Amiri font
%    - Company letterhead placeholder (name, address, contact)
%    - Reference/subject line (الموضوع)
%    - Recipient block and formal salutation
%    - Two-to-three-paragraph business body
%    - Closing, signature, cc list, and enclosures list
% ============================================================
\documentclass[12pt,a4paper]{article}

% --- Page geometry ---
\usepackage[a4paper,margin=2.5cm,top=2.5cm,bottom=2.5cm,headheight=16pt]{geometry}

% --- Standard packages (load ALL before bidi) ---
\usepackage{graphicx}
\usepackage{array}
\usepackage{booktabs}
\usepackage{tabularx}
\usepackage{enumitem}
\usepackage{titlesec}
\usepackage{fancyhdr}
\usepackage{setspace}
\usepackage{xcolor}
\usepackage{parskip}

% --- BIDI (must come before polyglossia, after standard packages) ---
\usepackage{bidi}

% --- Font specification (requires bidi) ---
\usepackage[no-math]{fontspec}

% --- Polyglossia (language support) ---
\usepackage[quiet]{polyglossia}

% --- Fonts ---
\setmainfont[Scale=1]{NotoKufiArabic-Regular.ttf}
\newfontfamily\arabicfont[Script=Arabic,Scale=0.95]{NotoKufiArabic-Regular.ttf}
\newfontfamily\englishfont[Scale=0.95]{TeX Gyre Termes}

% --- Languages ---
\setdefaultlanguage[calendar=gregorian,hijricorrection=1]{arabic}
\setotherlanguage[variant=british]{english}

% --- Hyperref (ALWAYS last) ---
\usepackage[breaklinks,hidelinks]{hyperref}
\hypersetup{pdftitle={خطاب تجاري},pdfauthor={اسم الشركة}}

% --- Line spacing ---
\setstretch{1.4}
\setlength{\parskip}{0.6em}
\setlength{\parindent}{0pt}

% --- Header/Footer ---
\pagestyle{fancy}
\fancyhf{}
\fancyhead[R]{\small\itshape خطاب تجاري}
\fancyhead[L]{\small اسم الشركة}
\fancyfoot[C]{\small\thepage}
\renewcommand{\headrulewidth}{0.3pt}

% ============================================================
%  Document
% ============================================================
\begin{document}

% --- Company letterhead placeholder ---
\begin{center}
{\LARGE\bfseries اسم الشركة}\\[0.3em]
% TODO: استبدل باسم الشركة الكامل
{\small العنوان: المدينة، الحي، الشارع، الرمز البريدي}\\[0.2em]
% TODO: استبدل ببيانات التواصل
{\small الهاتف: \LR{+000-000-000-0000} \,\textbar\, الفاكس: \LR{+000-000-000-0001} \,\textbar\, البريد: \LR{info@company.example}}
\end{center}

\vspace{0.5em}
\hrule
\vspace{1.5em}

% --- Reference block (DIN 5008-style: date + reference, right-aligned in RTL) ---
\begin{flushright}
% TODO: حدّث التاريخ
\textbf{التاريخ:} \today\\[0.2em]
% TODO: استبدل برقم المرجع الداخلي
\textbf{المرجع:} \LR{REF-2026-0001}
\end{flushright}

\vspace{1em}

% --- Recipient block ---
\begin{flushright}
% TODO: اسم المُرسَل إليه
السيد/ة \textbf{اسم المرسل إليه} المحترم/ة\\[0.2em]
% TODO: المسمى الوظيفي للمُرسَل إليه
المسمى: \textbf{المسمى الوظيفي}\\[0.2em]
% TODO: اسم الشركة المُرسَل إليها
\textbf{اسم الشركة المُرسَل إليها}\\[0.2em]
% TODO: عنوان المُرسَل إليه
العنوان: المدينة، الحي، الشارع، الرمز البريدي
\end{flushright}

\vspace{1em}

% --- Subject / reference line ---
\begin{flushright}
% TODO: حدّد موضوع الخطاب
\textbf{الموضوع:} \textit{عنوان الموضوع / موضوع الخطاب}
\end{flushright}

\vspace{0.5em}

% --- Salutation ---
\begin{flushright}
السيد/ة المحترم/ة،
\end{flushright}

\vspace{0.5em}

% --- Body: Paragraph 1 — Purpose / context ---
% TODO: اذكر الغرض من الخطاب والسياق
تحية طيبة وبعد،

يشرفني أن أتوجّه إليكم بهذا الخطاب نيابةً عن \textbf{اسم الشركة} بخصوص \textbf{موضوع الخطاب: عرض تجاري / استفسار / متابعة اتفاقية / دعوة تعاون}. يأتي هذا الخطاب استكمالاً لـ\textbf{المرجعية السابقة: اجتماع / مكالمة / مراسلة} بتاريخ \textbf{تاريخ المرجعية}، بهدف توضيح النقاط المتفق عليها والمضيّ قدماً نحو \textbf{الهدف التجاري}.

% --- Body: Paragraph 2 — Details / terms ---
% TODO: اذكر التفاصيل والشروط والبنود الرئيسية
نُقدّم فيما يلي أبرز النقاط: أولاً، \textbf{البند الأول: المنتج / الخدمة / العرض} بشروط \textbf{الشروط} ومدة تنفيذ \textbf{مدة التنفيذ}. ثانياً، \textbf{البند الثاني: السعر / الدفع / التسليم} وفقاً لـ\textbf{آلية الدفع المتفق عليها}. ثالثاً، \textbf{البند الثالث: الضمانات / الالتزامات} التي تضمن حقوق الطرفين. ونؤكد استعدادنا للالتزام بجميع البنود المذكورة أعلاه.

% --- Body: Paragraph 3 — Call to action / next steps ---
% TODO: اذكر الخطوة التالية والموعد المقترح
ونظراً لأهمية الموضوع، نرجو الردّ على هذا الخطاب في موعدٍ أقصاه \textbf{تاريخ الاستجابة المقترح}، مع تحديد \textbf{الخطوة التالية: اجتماع / توقيع عقد / زيارة ميدانية}. ويسعدنا تخصيص أي وقتٍ مناسب لمناقشة التفاصيل بشكلٍ مفصّل. وفي حال وجود أي استفسار، يُرجى التواصل معنا عبر القنوات الموضّحة في ترويسة الخطاب.

% --- Closing ---
\vspace{0.5em}
\begin{flushright}
وتفضلوا بقبول فائق الاحترام والتقدير،
\end{flushright}

\vspace{2em}

% --- Signature block ---
\begin{flushright}
% TODO: التوقيع اليدوي أعلى الاسم
\rule{5cm}{0.4pt}\\[0.3em]
\textbf{اسم المرسل}\\[0.2em]
% TODO: المسمى الوظيفي
المسمى الوظيفي\\[0.2em]
% TODO: اسم الشركة
اسم الشركة
\end{flushright}

\vspace{1.5em}
\hrule
\vspace{1em}

% --- CC list (نسخة لـ) ---
\begin{flushright}
\textbf{نسخة إلى (CC):}\\[0.2em]
% TODO: أضف أسماء الجهات المُرسَل إليها نسخة
\begin{itemize}[leftmargin=2em,rightmargin=0pt,nosep,itemsep=0.2em]
  \item \textbf{اسم الجهة الأولى} --- المسمى الوظيفي
  \item \textbf{اسم الجهة الثانية} --- المسمى الوظيفي
  \item \textbf{اسم الجهة الثالثة} --- المسمى الوظيفي
\end{itemize}
\end{flushright}

\vspace{0.5em}

% --- Enclosures list (المرفقات) ---
\begin{flushright}
\textbf{المرفقات (Enclosures):}\\[0.2em]
% TODO: أضف قائمة بالمرفقات
\begin{itemize}[leftmargin=2em,rightmargin=0pt,nosep,itemsep=0.2em]
  \item \textbf{المرفق الأول} --- وصف موجز
  \item \textbf{المرفق الثاني} --- وصف موجز
  \item \textbf{المرفق الثالث} --- وصف موجز
\end{itemize}
\end{flushright}

\end{document}
